\section{Creation, Grace, and Glory}\label{sec:creation}

Having treated the angels' nature, knowledge, and will, Aquinas turns to
``their creation, or their origin taken universally'' (\latin{de eorum
creatione, sive universaliter de eorum exordio}). The inquiry is
threefold: how the angels were produced in natural being (q.~61), how
they were perfected in grace and glory (q.~62), and how some of them
became evil (qq.~63--64, \S\ref{sec:fall}) (\ST{I q.~61 pr.}). The
Catechism gathers the first two moments into the Church's confession:
the angels ``were created through and for'' Christ (Col 1:16), and they
``have been present since creation and throughout the history of
salvation'' (CCC~331--332).

\subsection{The Production of the Angels in the Order of Natural Being
(\ST{I q.~61})}\label{sec:q61}

The four articles ask whether the angel has a cause of his being,
whether he was produced from eternity, whether he was created before the
corporeal world, and whether the angels were created in the empyrean
heaven (\ST{I q.~61 pr.}).

\paragraph{Created by God (a.~1).} The first objection observes that
Genesis never mentions the angels; the \emph{sed contra} answers with
Ps 148:2 and~5: ``Praise ye him, all his angels, praise ye him, all his
hosts \dots\ For he spoke, and they were made'' (Douay). ``It must be
affirmed that angels and everything existing, except God, were made by
God'': God alone is his own being; everything else has being by
participation, and ``whatever exists by participation is caused by what
exists essentially'' (\latin{omnia alia sunt entia per participationem}).
Augustine had already answered the silence of Genesis: the angels ``are
that light which was called `Day,'\,'' for ``if we are justified in
understanding in this light the creation of the angels, then certainly
they were created partakers of the eternal light which is the
unchangeable Wisdom of God'' --- in his Latin, \latin{ut ea luce
inluminati, qua creati, fierent lux et uocarentur dies} (\emph{De
civitate Dei} XI.9). Moses, Aquinas adds, was ``addressing an uncultured
people, as yet incapable of understanding an incorporeal nature,'' and
to name spirits above all bodies would have been an occasion of idolatry
(ad~1). The Church has since defined the substance of the article:
Lateran~IV confesses that God ``by his almighty power at the beginning
of time created at once out of nothing both creatures, the spiritual and
the corporeal, namely the angelic and the earthly'' --- \latin{simul ab
initio temporis utramque de nihilo condidit creaturam, spiritualem et
corporalem, angelicam videlicet et mundanam} (\emph{Firmiter}, DS~800) ---
and Vatican~I, \emph{Dei Filius} ch.~1, repeats the same words of the
angels' creation \latin{de nihilo}.

\paragraph{Not from eternity (a.~2).} The \emph{sed contra} is
Prov 8:22 in the person of begotten Wisdom, who was possessed ``before
He made anything from the beginning.'' ``God alone, Father, Son and Holy
Ghost, is from eternity. Catholic Faith holds this without doubt; and
everything to the contrary must be rejected as heretical.'' To be
created \latin{ex nihilo} is to be made ``after'' not being: \latin{idest
postquam nihil fuerat}. The replies guard the angel's lofty duration
without yielding the point. The divine will, not a necessity of the
divine being, produced creatures ``such as He willed, and when He
willed'' (ad~1). The angel is above the time that measures the heavens,
but not above the succession of his own being after non-being; Aquinas
cites Augustine (\emph{De Genesi ad litteram} VIII.20--21) that ``God
moves the spiritual creature according to time'' (ad~2). The capacity
for truth that makes the angelic nature incorruptible was bestowed when
God willed, not possessed from eternity (ad~3). The angel's own measure
of duration, the \latin{aevum}, is treated in the commentary on the
Sentences as a participated eternity between time and God's
(\emph{In II Sent.} d.~2 q.~1 a.~1).

\paragraph{Created with the corporeal world, not before it (a.~3).} The
objections carry the Greek line. Jerome, commenting on Titus 1:2, asks
--- in Aquinas's citation --- ``Six thousand years of our time have not
yet elapsed; yet how shall we measure the time, how shall we count the
ages, in which the Angels, Thrones, Dominations, and the other orders
served God?'' Jerome's own text runs: \latin{in quibus angeli, throni,
dominationes, caeteraeque virtutes servierint Deo, et absque temporum
vicibus atque mensuris, Deo jubente, substiterint} (\emph{In ep.~ad
Titum} 1:2, PL~26). Damascene reports the same opinion: ``Some, indeed,
like Gregory the Theologian, say that these were before the creation of
other things,'' and for himself, ``it was fitting that the mental
essence should be the first created, and then that which can be
perceived, and finally man'' (\emph{De fide orthodoxa} II.3). Gregory's
own oration says: ``He first conceived the Heavenly and Angelic Powers.
And this conception was a work fulfilled by His Word, and perfected by
His Spirit''; then, ``when His first creation was in good order, He
conceives a second world, material and visible'' (\emph{Oration} 38.9--10).
Aquinas answers: ``There is a twofold opinion on this point to be found
in the writings of the Fathers. The more probable one holds that the
angels were created at the same time as corporeal creatures.'' The
angels ``are part of the universe: they do not constitute a universe of
themselves''; ``no part is perfect if separate from the whole''; and it
is improbable that God, whose works are perfect (Deut 32:4), created the
angelic creature apart, before the rest. Yet ``the contrary is not to be
deemed erroneous; especially on account of the opinion of Gregory
Nazianzen, `whose authority in Christian doctrine is of such weight that
no one has ever raised objection to his teaching, as is also the case
with the doctrine of Athanasius,' as Jerome says.'' Jerome, he notes,
spoke ``according to the teaching of the Greek Fathers; all of whom hold
the creation of the angels to have taken place previously to that of the
corporeal world'' (ad~1). On the alternative reading, ``In the
beginning'' of Gen 1:1 means ``in the Son'' or ``in the beginning of
time,'' not ``before which there was nothing'' simply. The same openness
marks the parallel treatments: in the Sentences commentary Aquinas calls
simultaneous creation more probable ``to exclude the error'' of angels
as creators of bodies, ``yet the other sentence is not to be prejudged
\dots\ because it is not demonstrated, nor expressed by faith''
(\latin{non est demonstratum, nec fide expressum}; \emph{In II Sent.}
d.~2 q.~1 a.~3), and in \emph{De potentia} the prior-creation opinion,
``held by great doctors, namely Basil, Gregory Nazianzen, and certain
others, is not to be reprobated as erroneous,'' though he judges
Augustine's opinion ``more reasonable'': \latin{Angeli simul cum creatura
corporali sunt conditi; tamen sine alterius opinionis praeiudicio}
(\emph{De potentia} q.~3 a.~18). Lateran~IV's \latin{simul ab initio
temporis} (\emph{Firmiter}, DS~800) professes the simultaneity of both
creatures' production at time's beginning; Aquinas, writing after the
council, still left the precedence question to the schools, and this
document leaves it where he left it.

\paragraph{In the empyrean heaven (a.~4).} The \emph{sed contra} is
Strabus on Gen 1:1: ``By heaven he does not mean the visible firmament,
but the empyrean, that is, the fiery or intellectual firmament, which is
not so styled from its heat, but from its splendor; and which was filled
with angels directly it was made'' (cited by Aquinas). Since corporeal
and spiritual creatures constitute one universe, and spiritual creatures
were created ``to bear some relationship to the corporeal creature, and
to rule over every corporeal creature,'' it was fitting ``for the angels
to be created in the highest corporeal place, as presiding over all
corporeal nature; whether it be styled the empyrean heaven
(\latin{caelum empyreum}), or whatever else it be called''; so Isidore,
cited by Aquinas, calls the highest heaven ``the heaven of the angels''
on Deut 10:14. Their creation in a place implies no dependence on a
body, for existence or for making --- ``God could have created them
before all corporeal creation, as many holy Doctors hold'' --- but shows
their order to corporeal nature, ``and that they are by their power in
touch with bodies'' (ad~1). Augustine, whom the second objection cites
(\emph{De Genesi ad litteram} III.10) as placing the angels' creation in
the upper atmosphere, perhaps means the highest heaven by that name, or
speaks only of the angels that sinned (ad~2). The ``heaven'' of
Isa 14:13, into which the sinning angel boasted to ascend, is no
corporeal heaven but ``the heaven of the Blessed Trinity'' (ad~3).

\angeldossier{\ST{I q.~61 aa.~1--4}; parallels \emph{In II Sent.} d.~2
q.~1 a.~3, d.~2 q.~2; \emph{De potentia} q.~3 aa.~17--18.}
 {Defined that the angels were created by God out of nothing at the
 beginning of time, together with corporeal creation as such (Lateran
 IV, \emph{Firmiter}, DS~800: \latin{simul ab initio temporis utramque
 de nihilo condidit creaturam, spiritualem et corporalem, angelicam
 videlicet et mundanam}; renewed by Vatican~I, \emph{Dei Filius} ch.~1),
 and that the angelic creature had a beginning and is not eternal (a.~2;
 Aquinas: the contrary ``must be rejected as heretical''). Church
 teaching that the angels were created through and for Christ and have
 been present since creation (CCC~331--332). Disputed whether the angels
 were created before the corporeal world (a.~3): Aquinas holds
 simultaneous creation the more probable opinion yet refuses to brand
 the Greek view erroneous; he records the same restraint in \emph{In II
 Sent.} d.~2 q.~1 a.~3 (\latin{nec fide expressum}) and \emph{De
 potentia} q.~3 a.~18. Common teaching of the Latin Doctors, as a
 judgment of fitness and not of faith, that the angels were created in
 the empyrean or highest heaven (a.~4; Strabus and Isidore, cited by
 Aquinas).}
 {Ps 148:2,~5 (\emph{sed contra} of a.~1); Prov 8:22 (\emph{sed contra}
 of a.~2); Gen 1:1 and Deut 32:4 (a.~3); Isa 14:13 and Deut 10:14
 (a.~4); Augustine, \emph{De civitate Dei} XI.9 (a.~1 ad~1, read at the
 locus), \emph{De Genesi ad litteram} VIII.20--21 and III.10 (cited by
 Aquinas); Jerome, \emph{In ep.~ad Titum} 1:2 (a.~3 obj.~1, read at
 PL~26); Damascene, \emph{De fide orthodoxa} II.3 (a.~3 obj.~1, read at
 the locus); Gregory Nazianzen, \emph{Oration} 38.9--10 (read at the
 locus); Strabus, \emph{Glossa} on Gen 1:1, and Isidore on Deut 10:14
 (\emph{sed contra} and corpus of a.~4, cited by Aquinas); Lateran~IV,
 \emph{Firmiter} (DS~800); Vatican~I, \emph{Dei Filius} ch.~1;
 CCC~331--332.}
 {The Greek Doctors --- Gregory Nazianzen (\emph{Or.} 38.9--10),
 Damascene (\emph{De fide orth.} II.3), and Jerome reporting the Greeks
 --- place the angels' creation before the corporeal world (a.~3);
 Augustine places it ``in the upper atmosphere'' as the second objection
 reads him (a.~4 obj.~2, cited by Aquinas).}

\subsection{The Perfection of the Angels in the Order of Grace and of
Glory (\ST{I q.~62})}\label{sec:q62}

Nine articles: whether the angels were created in beatitude; whether
they needed grace to turn to God; whether they were created in grace;
whether they merited beatitude; whether they obtained it immediately
after one meritorious act; whether grace and glory were given according
to the degree of their natural gifts; whether natural knowledge and love
remain in the blessed; whether a beatified angel can sin; and whether
the blessed angels advance in beatitude (\ST{I q.~62 pr.}).

\paragraph{Not created in beatitude (a.~1).} ``To be established or
confirmed in good is of the nature of beatitude. But the angels were not
confirmed in good as soon as they were created; the fall of some of them
shows this'' (\emph{sed contra}). Beatitude names the ultimate
perfection of an intellectual nature, and that perfection is twofold:
one the angel can reach by his natural power, ``in a measure called
beatitude or happiness''; above it stands the beatitude by which ``we
shall see God as He is,'' which ``is beyond the nature of every created
intellect'' (\latin{beatitudo supernaturalis}). ``As regards this first
beatitude, which the angel could procure by his natural power, he was
created already blessed,'' since he receives his natural perfection
without discourse; ``but the angels did not have from the beginning of
their creation that ultimate beatitude which is beyond the power of
nature; because such beatitude is no part of their nature, but its
end.'' The first objection's text from \emph{De ecclesiasticis
dogmatibus} speaks of that natural perfection in the state of innocence
(ad~1); and Augustine's ``morning knowledge'' --- the angel's knowledge
of things in the Word --- was present from creation only in its natural
mode; the knowledge of glory came ``when the angels became blessed by
turning to the good'' (ad~3, citing \emph{De Genesi ad litteram} IV.34,
V.5).

\paragraph{The need of grace to turn to God (a.~2).} The \emph{sed
contra} is Rom 6:23: ``The grace of God is life everlasting.'' The
will's natural inclination reaches only what accords with its nature;
toward what is above nature it cannot be inclined ``unless helped by
some other supernatural principle.'' Since to see God in his essence,
``wherein the ultimate beatitude of the rational creature consists, is
beyond the nature of every created intellect, \dots\ no rational
creature can have the movement of the will directed towards such
beatitude, except it be moved thereto by a supernatural agent. This is
what we call the help of grace'' (\latin{auxilium gratiae}). The angel's
natural love of God, treated above (a.~2 ad~1; \S\ref{sec:q60}), has God
as ``the author of his natural being'' for its object, not God as
beatifying by the vision of his essence. For man the conversion is
difficult both as supernatural and as hindered by flesh and sin; for the
angel, ``only because it is supernatural'' (ad~2). There is no infinite
regress: conversion by perfect love requires consummate grace
(\latin{gratia consummata}), meritorious conversion requires habitual
grace, and the conversion that prepares for grace requires only God's
operation drawing the soul --- ``Convert us, O Lord, to Thee, and we
shall be converted'' (Lam 5:21, ad~3).

\paragraph{Created in grace (a.~3).} The \emph{sed contra} is
Augustine's question and answer: ``who made this will, which already
they had, but He who created them with a good will, or with that chaste
love by which they cleaved to Him, in one and the same act creating
their nature, and endowing it with grace?'' (\emph{De civitate Dei}
XII.9) --- \latin{simul eis et condens naturam et largiens gratiam}.
``Although there are conflicting opinions on this point, some holding
that the angels were created only in a natural state, while others
maintain that they were created in grace; yet it seems more probable,
and more in keeping with the sayings of holy men, that they were created
in sanctifying grace'' (\latin{in gratia gratum faciente}). The argument
is by Augustine's seedlike forms: as the \latin{rationes seminales} of
all natural effects were implanted in the first creation, so grace,
which bears to beatitude the relation of seed to effect --- ``grace is
called the `seed' of God'' (1 John 3:9) --- was given with the angelic
nature from the beginning. The Master's school stood otherwise. The
Lombard teaches that the angels at creation were good and innocent but
``could not advance to the merit of life unless grace were superadded,
which was added to certain ones in their confirmation'' (\latin{non
poterant proficere ad meritum vitae, nisi gratia superadderetur, quae
addita est quibusdam in confirmatione}; \emph{Sent.} II d.~3 c.~4), and
that the fallen were deserted by a grace ``not so that it was withdrawn
after being given, but because it was never applied to turn them''
(\latin{non ita quod prius dedita subtraheretur, sed quia nunquam est
apposita ut converterentur}; \emph{Sent.} II d.~5 c.~1): ``cooperating
grace was given to the angels who stood, by which they were converted''
(\emph{Sent.} II d.~5 c.~4). Aquinas's own Sentences commentary records
that the natural-state opinion ``is the more common'' (\latin{haec
opinio est communior}; \emph{In II Sent.} d.~4 q.~1 a.~3); in the
\emph{Summa} he judges creation in grace the more probable and the more
consonant with the saints. Grace, the replies add, does not force the
intellectual nature it inclines: ``he who has grace can fail to make use
of it, and can sin'' (ad~2); and grace, as principle of right operation
rather than end, was fittingly given at once with nature (ad~3).

\paragraph{Beatitude merited (a.~4).} The \emph{sed contra} is
Apoc 21:17, where the heavenly Jerusalem has ``the measure of a man,
that is, of an angel'': as man reaches beatitude by merit, so does the
angel. ``Perfect beatitude is natural only to God, because existence and
beatitude are one and the same thing in Him. Beatitude, however, is not
of the nature of the creature, but is its end.'' Every operation that
attains an end beyond the agent's power is not productive but deserving
--- \latin{meritoria} --- and since ultimate beatitude exceeds both
angelic and human nature, ``it remains, then, that both man and angel
merited their beatitude.'' The alternatives collapse: beatitude without
grace and merit would contradict the very notion of an end and reward;
meriting within beatitude contradicts the notion of \latin{meritum}, for
what is at its term is not moved toward it; and one and the same act
cannot be informed by imperfect grace, the principle of meriting, and by
perfect grace, the principle of enjoying. ``Consequently it is better to
say that the angel had grace ere he was admitted to beatitude, and that
by such grace he merited beatitude.'' The angel's difficulty lay not in
contrary inclinations but in the work's exceeding his natural capacity
(ad~1); and he merited ``by the movement of charity, which comes of
grace,'' not by the natural movement toward God (ad~2).

\paragraph{Beatitude after one meritorious act (a.~5).} Man's soul and
the angel are alike ordained to beatitude, and the separated soul with
merit enters beatitude at once unless something obstructs; the angel,
who had merit in his first act of charity and no obstacle, ``passed at
once into beatitude by only one meritorious act'' (\emph{sed contra}).
``The angel was beatified instantly after the first act of charity,
whereby he merited beatitude,'' for ``grace perfects nature according to
the manner of the nature'' (\latin{gratia perficit naturam secundum
modum naturae}), and it belongs to the angelic nature to possess its
perfection not by stages but at once. Even in man every act informed by
charity merits beatitude; much more did the angel's one act suffice. A
longer way was assigned to man, who was not intended to reach his
perfection at once (ad~1). The angel is above corporeal time, so his
``instants'' are counted by the succession of his acts: one instant for
the meritorious act, another for the beatific fruition, since imperfect
and consummate grace cannot inform one act together (ad~2).

\paragraph{Grace and glory according to natural degree (a.~6).} The
\emph{sed contra} is the Master of the Sentences: ``those angels who
were created with more subtle natures and of keener intelligence in
wisdom, were likewise endowed with greater gifts of grace'' (Peter
Lombard, \emph{Sent.} II d.~3 c.~2). ``It is reasonable to suppose that
gifts of graces and perfection of beatitude were bestowed on the angels
according to the degree of their natural gifts,'' for two reasons. On
God's part, the grades of nature were made for the grades of grace and
glory, as the builder who chisels some stones more finely sets them
apart for the more honored place. On the angel's part, nothing in him
thwarts or retards the intellectual movement, as the sensitive appetite
does in man; ``when there is nothing to retard or thwart it, nature is
moved with its whole energy,'' so the higher natures ``were turned to
God more mightily and efficaciously.'' God's will, which is the sole
source of grace, is also the source of nature and of its degrees (ad~1);
and the argument does not transfer to men, who differ only numerically
and whose intellectual movement a lower appetite can impede (ad~3).

\paragraph{Natural knowledge and love remain (a.~7).} The \emph{sed
contra}: ``So long as a nature endures, its operation remains. But
beatitude does not destroy nature, since it is its perfection.'' The
principles of operations stand to one another as the operations
themselves; nature is to beatitude as first to second, ``but the first
must ever be preserved in the second. Consequently nature must be
preserved in beatitude: and in like manner the act of nature must be
preserved in the act of beatitude'' (\latin{semper oportet salvari
primum in secundo}). The perfection removes the opposed privation, not
the underlying imperfection of nature: the angel can know God by his
essence --- the knowledge of glory --- and at the same time know God by
his own essence, which belongs to his natural knowledge, as one may know
the same thing through a probable and through a demonstrative medium
(ad~1). The gifts of glory presuppose the natural gifts, since no
beatitude is self-subsisting save the uncreated (ad~2); and two acts of
one faculty can stand together when one is ordained to the other, as
natural knowledge and love are ordained to the knowledge and love of
glory (ad~3).

\paragraph{The blessed cannot sin (a.~8).} Aquinas cites Augustine
(\emph{De Genesi ad litteram} XI): ``there is in the holy angels that
nature which cannot sin'' (\emph{sed contra}). ``The beatified angels
cannot sin,'' because their beatitude consists in seeing God through his
essence, and God's essence is the very essence of goodness; the angel
beholding God is therefore disposed to him as any will is disposed to
the common aspect of good --- it cannot will or act except aiming at
good --- ``therefore the beatified angel can neither will nor act,
except as aiming towards God,'' and so cannot sin. Created good,
considered in itself, can fail; but its perfect union with the uncreated
good in beatitude renders it unable to sin (ad~1). The will is referred
to opposites in what lies beneath it, not in what it naturally tends to;
toward God, whom the blessed see to be the very nature of goodness,
there is no alternative aim (ad~2). Free choice stands to the means as
the intellect to conclusions: liberty is perfected by ranging among the
means within the order of the end, and to choose by turning from that
order --- which is sin --- is a defect of liberty, not its exercise.
``Hence there is greater liberty of will in the angels, who cannot sin,
than there is in ourselves, who can sin'' (ad~3).

\paragraph{No advance in beatitude (a.~9).} Merit and progress belong to
the state of the way; the angels are not \latin{viatores} but
\latin{comprehensores} (\emph{sed contra}). Every movement is directed
by the mover's intention to one determinate end, and infinite progress
is repugnant to an end. The rational creature is led by God to the
vision of the supreme Essence, not to the supreme mode of vision, which
is comprehension and belongs to God alone; between finite efficacy and
the infinite object the degrees of vision are without number, and each
creature is brought by predestination to its determinate degree.
``Consequently, when that degree is once secured, it cannot pass to a
higher degree.'' Perfect charity does not merit but enjoys the reward,
as the act of an acquired habit no longer produces the habit (ad~1).
The angels' ministries profit them not as merits toward an end but as
parts of their beatitude: ``to pour out acquired perfection upon others
is of the nature of what is perfect, considered as perfect'' (ad~2).
Their joy can still grow at the salvation of those saved by their
ministry --- ``There is joy before the angels of God upon one sinner
doing penance'' (Luke 15:10) --- but this belongs to the accidental
reward, which can increase until judgment day. Some writers say the
blessed can merit in that respect; Aquinas judges ``it is better to say
that the Blessed can in no wise merit without being at the same time a
wayfarer and a comprehensor; like Christ, Who alone was such'' (ad~3).

\angeldossier{\ST{I q.~62 aa.~1--9}; parallels \emph{In II Sent.} d.~4
q.~1 aa.~1--3.}
 {Common teaching that the angels needed grace to turn to God as the
 object of beatitude (a.~2), that they merited their beatitude (a.~4;
 Apoc 21:17) and obtained it immediately upon their first act of charity
 (a.~5), that natural knowledge and love remain in the blessed (a.~7),
 and that the confirmed angels can neither sin (a.~8) nor advance in
 essential beatitude (a.~9). Common teaching, with the Master of the
 Sentences, that grace and glory were given according to the degrees of
 nature (a.~6; Peter Lombard, \emph{Sent.} II d.~3 c.~2). Disputed
 whether the angels were created in sanctifying grace (a.~3): Aquinas
 holds it the more probable opinion, consonant with Augustine (\emph{De
 civitate Dei} XII.9); the Lombard (\emph{Sent.} II dd.~3, 5) and, as
 Aquinas's commentary attests, the more common opinion of the schools
 hold grace to have been superadded after creation. Disputed likewise,
 in a lesser register, whether the blessed merit an increase of
 accidental reward (a.~9 ad~3); Aquinas prefers the denial.}
 {Augustine, \emph{De civitate Dei} XII.9 (\emph{sed contra} of a.~3,
 read at the locus), \emph{De Genesi ad litteram} IV.34, V.5 (a.~1
 obj.~3), XI (\emph{sed contra} of a.~8), cited by Aquinas;
 \emph{De ecclesiasticis dogmatibus} (a.~1 obj.~1, cited by Aquinas);
 Peter Lombard, \emph{Sent.} II d.~3 cc.~2,~4 and d.~5 cc.~1,~4
 (read at the locus); Rom 6:23 (\emph{sed contra} of a.~2); 1 John 3:9
 (a.~3); Apoc 21:17 (\emph{sed contra} of a.~4); 1 Cor 13:10 (a.~7
 obj.~1); Luke 15:10 (a.~9 ad~3); Lam 5:21 (a.~2 ad~3).}
 {Peter Lombard, \emph{Sent.} II d.~3 c.~4 and d.~5 cc.~1,~4: the angels
 created without grace, cooperating grace given afterward to those who
 stood (a.~3); the ``more common'' school opinion to the same effect
 recorded in \emph{In II Sent.} d.~4 q.~1 a.~3; the writers who let the
 blessed merit their accidental reward (a.~9 ad~3).}
